\section{The radial inequality}\label{rad:section}

This appendix proves the pointwise inequality used for positive channel
parameters.  Its algebraic certificates are given explicitly, so the proof
requires no numerical or computer-assisted verification.

\begin{lemma}[Radial inequality]\label{rad:lemma}
Let $m=(m_1,m_2,m_3)$ have positive coordinates and $\av{m}=1$, where
$\av{w}=\frac13\sum_{i=1}^3w_i$.  Put
\[
 v=m-1,\qquad x=\frac12\av{v^2},\qquad
 g=\av{\log m},\qquad u_i=\log m_i-g,
\]
\[
 J=\frac12\av{v\log m},\qquad B=\frac12\av{m u^2},
 \qquad F(m)=3g+4J+\frac45B.
\]
Since $I=g+2J$, this is the functional $F$ in
Equation~\eqref{pos:functional-F}. Then, for every $0<\lambda<1$,
\begin{equation}\label{rad:inequality}
 F(m)-\lambda^{-2}F(1+\lambda v)
 \geq \frac{1-\lambda^2}{2}\,x^2.
\end{equation}
\end{lemma}

\begin{proof}
Fix the message in the statement and
write
\[
 m(\tau)=1+\tau v^{(0)},\qquad x(\tau)=\tau^2x_0,
 \qquad \mathcal F(\tau)=F(m(\tau)),\qquad 0\leq\tau\leq1,
\]
where $v^{(0)}=m-1$ and $x_0=\frac12\av{(v^{(0)})^2}$.
Every coordinate of $m(\tau)$ is positive on this interval.
The functions are analytic near $\tau=0$.  Since $\av{v^{(0)}}=0$,
we have $g(\tau)=O(\tau^2)$, $J(\tau)=O(\tau^2)$ and
$B(\tau)=O(\tau^2)$, so $\mathcal F(\tau)=O(\tau^2)$.
Let $\partial=\tau\frac{d}{d\tau}$. We claim that it suffices to prove
\begin{equation}\label{rad:differential-bound}
 \partial(\partial-1)(\partial-2)\mathcal F(\tau)
 \geq12x(\tau)^2\qquad(0<\tau\leq1).
\end{equation}
To see this, suppose the bound holds along the ray.
Define
\[
 Y=(\partial-2)\mathcal F-x(\tau)^2,\qquad Z=(\partial-1)Y.
\]
The operator $\partial-2$ annihilates the quadratic Taylor term, giving
$Y=O(\tau^3)$ and $Z=O(\tau^3)$ as $\tau\downarrow0$.
Since $\partial(\partial-1)x(\tau)^2=12x(\tau)^2$,
Equation~\eqref{rad:differential-bound} gives $\partial Z\geq0$.
For $\tau>0$ this means $Z'(\tau)\geq0$, and $Z(0)=0$ therefore gives
$Z\geq0$.  Next,
\[
 \left(\frac{Y(\tau)}{\tau}\right)'
   =\frac{Z(\tau)}{\tau^2}\geq0.
\]
The limit of $Y(\tau)/\tau$ at zero is zero, so $Y\geq0$.  Finally,
\[
 \frac{d}{d\tau}
 \left[\frac{\mathcal F(\tau)-\tfrac12\tau^4x_0^2}{\tau^2}\right]
 =\frac{Y(\tau)}{\tau^3}\geq0.
\]
Evaluating the resulting monotonicity between $\tau=\lambda$ and $\tau=1$
gives
\[
 F(m)-\frac12x_0^2
 \geq\lambda^{-2}F(1+\lambda v^{(0)})
          -\frac12\lambda^2x_0^2,
\]
which is Equation~\eqref{rad:inequality}.

It remains to establish Equation~\eqref{rad:differential-bound} for every
positive message. All quantities in the following calculation refer to the
current message on the ray. Set
\[
 y=\frac12\av{v^3}.
\]
Thus $\partial v=v$, $\partial x=2x$, and $\partial y=3y$.
Introduce
\[
 U_j=\av{m^{-j}},\qquad S_j=\av{u m^{-j}},\qquad
 U=U_1,\quad S=S_1,\quad M=2J,\quad V=2B,\quad W=\av{u^2}.
\]
In particular $M=\av{m u}$, since $\av{u}=0$, and
$F=3g+2M+\frac25V$.  The identities
\[
 \partial\log m_i=1-m_i^{-1},\qquad
 \partial u_i=U-m_i^{-1}
\]
give
\begin{equation}\label{rad:differentiation}
\begin{aligned}
 \partial g&=1-U,&
 \partial M&=M+U-1,\\
 \partial V&=V-W+2UM,&
 \partial W&=-2S,\\
 \partial S&=U^2-U_2-S+S_2,&
 \partial U_j&=j(U_{j+1}-U_j).
\end{aligned}
\end{equation}
For example, differentiating $V=\av{m u^2}$ gives
$\av{v u^2}+2\av{m u(U-m^{-1})}=V-W+2UM$;
differentiating $S$ gives the next-to-last identity in
\eqref{rad:differentiation}.  Applying these identities to $F$ yields
\begin{equation}\label{rad:G}
\begin{aligned}
 G&:=\partial(\partial-1)(\partial-2)F\\
  &=6-12U+8U_2-2U_3\\
  &\quad+\frac45\bigl[S_2-3S+(2U_3-3U_2)M
                    +3(U-1)(U_2-U)\bigr].
\end{aligned}
\end{equation}
Our remaining goal is $G\geq12x^2$.

The elementary symmetric identities for the three zero-sum coordinates $v_i$
are
\begin{equation}\label{rad:cubic}
 v_i^3=3xv_i+2y,\qquad
 D_t:=\prod_{i=1}^3(1+t v_i)=1-3t^2x+2t^3y.
\end{equation}
Write $D=D_1$.  For $0\leq t\leq1$ all the factors in $D_t$ are positive
and have arithmetic mean one, so
\begin{equation}\label{rad:denominator-range}
 0<D_t\leq1.
\end{equation}
Define the coordinate polynomials
\[
 h_i=1-v_i+v_i^2-3x,\qquad
 k_i(t)=v_i-t(v_i^2-2x).
\]
The product of the two factors $1+v_j$ with $j\ne i$ is $h_i$, hence
$m_i^{-1}=h_i/D$.  Similarly, centering the logarithmic integral gives
\begin{equation}\label{rad:log-integral}
 \frac{v_i}{1+t v_i}-\av{\frac{v}{1+t v}}
   =\frac{k_i(t)}{D_t},\qquad
 u_i=\int_0^1\frac{k_i(t)}{D_t}\,dt.
\end{equation}
The first identity follows by multiplying by $D_t$ and using
\eqref{rad:cubic}; the second follows by integration.  All denominators here
are bounded away from zero for the fixed positive message under consideration.

Put $H_j=\av{h^j}$.  Substitution of these formulas into \eqref{rad:G}
separates its polynomial and logarithmic terms:
\begin{equation}\label{rad:RN-representation}
 D^3(G-12x^2)=R+\int_0^1\frac{N(t)}{D_t}\,dt,
\end{equation}
where
\begin{equation}\label{rad:RN-definition}
\begin{aligned}
 R={}&(6-12x^2)D^3-12H_1D^2+8H_2D-2H_3
       +\frac{12}{5}(H_1-D)(H_2-H_1D),\\
 N(t)={}&\frac45\av{k(t)
             \bigl[h^2D-3hD^2+(2H_3-3H_2D)v\bigr]}.
\end{aligned}
\end{equation}
In particular $N(t)$ is affine in $t$.  Let
\begin{equation}\label{rad:L-definition}
 L=R+\int_0^1N(t)(2-D_t)\,dt.
\end{equation}
Since
\[
 \frac1{D_t}-(2-D_t)=\frac{(1-D_t)^2}{D_t}\geq0,
\]
it suffices for \eqref{rad:differential-bound} to prove
\begin{equation}\label{rad:polynomial-goal}
 N(0)\geq0,\qquad N(1)\geq0,\qquad L\geq0.
\end{equation}
Indeed, the first two inequalities imply $N(t)\geq0$ on $[0,1]$, and then
\eqref{rad:RN-representation} is at least $L$.

We give the polynomial computations and their sign certificates in full.
For their verification one only needs the recurrence
\[
 P_0=1,\quad P_1=0,\quad P_2=2x,\qquad
 P_j=3xP_{j-2}+2yP_{j-3}\quad(j\geq3),
 \qquad P_j=\av{v^j}.
\]
The remaining moments that occur in \eqref{rad:RN-definition} are
\[
 P_3=2y,\quad P_4=6x^2,\quad P_5=10xy,\quad
 P_6=18x^3+4y^2.
\]
Expanding $h^j$ and using these moments gives
\begin{equation}\label{rad:H-expansions}
\begin{aligned}
 H_1&=1-x,\\
 H_2&=3x^2-4y+1,\\
 H_3&=-9x^3+9x^2+6xy+3x+4y^2-14y+1.
\end{aligned}
\end{equation}
The two coefficients of the affine polynomial $N(t)$ are explicitly
\begin{equation}\label{rad:N-expansion}
\begin{aligned}
 \frac5{24}N(t)={}&A(x,y)+t C(x,y),\\
 A(x,y)={}&3x^4+12x^3+x^2
   +(-6x^3-28x^2-2x)y+(26x-2)y^2-4y^3,\\
 C(x,y)={}&6x^4-2x^3
   +(-13x^3-10x^2-x)y+(10x^2+22x+4)y^2-16y^3.
\end{aligned}
\end{equation}
In particular their sum factors as
\begin{equation}\label{rad:N1}
 \frac5{24}N(1)=(x-y)(9x-10y+1)(x^2+x-2y).
\end{equation}
Integrating \eqref{rad:N-expansion} against $2-D_t$ in
\eqref{rad:L-definition}, and using \eqref{rad:H-expansions}, gives
\begin{equation}\label{rad:L-expansion}
\begin{aligned}
 \frac{25}{6}L={}&1500x^5-1020x^4+240x^3+168y^4\\
 &+(-480x^2-756x-52)y^3\\
 &+(2114x^3+710x^2+48x-80)y^2\\
 &+(-3093x^4+736x^3-35x^2)y.
\end{aligned}
\end{equation}
These are finite polynomial identities: the displayed moment recurrence
provides a direct way to check every coefficient.

We next describe the domain on which their signs are required.  Positivity
and $\sum_i m_i=3$ imply $\sum_i m_i^2<9$, and therefore $0\leq x<1$.
The discriminant identity for \eqref{rad:cubic} is
\[
 \prod_{1\leq i<j\leq3}(v_i-v_j)^2=108(x^3-y^2),
\]
so $|y|\leq x^{3/2}$.  Consequently each factor in \eqref{rad:N1} is
nonnegative, as is seen from
\begin{align*}
 x-y&\geq x(1-\sqrt{x}),\\
 9x-10y+1&\geq(1-\sqrt{x})(10x+\sqrt{x}+1),\\
 x^2+x-2y&\geq x(1-\sqrt{x})^2.
\end{align*}
This proves $N(1)\geq0$.

If $x=0$, then $v=0$ and the lemma is immediate, so henceforth take $x>0$.
Set $r=y/x$.  Then $r^2\leq x$ and $r\leq\sqrt{x}<1$, while $D>0$ gives
\begin{equation}\label{rad:domain}
 r^2\leq x<\frac1{3-2r}.
\end{equation}
The denominator is positive.  Combining these inequalities yields
\[
 0<1-r^2(3-2r)=(1-r)^2(2r+1),
\]
and hence $-\frac12<r<1$.  Every point of \eqref{rad:domain} is represented by
\begin{equation}\label{rad:z-substitution}
 x=\frac{r^2+z/(3-2r)}{1+z},\qquad y=rx,\qquad
 z=\frac{x-r^2}{1/(3-2r)-x}\geq0.
\end{equation}
Thus the parametrization includes $x=r^2$ without a limiting argument;
the excluded upper endpoint corresponds to a zero coordinate of $m$.

For clarity write $\widetilde N_0(r,z)$ and $\widetilde L(r,z)$ for $N(0)$
and $L$ after the substitution \eqref{rad:z-substitution}, and keep $x$
on the left sides below as the expression in that substitution.  The exact
coefficient identity for $N(0)$ is
\begin{equation}\label{rad:N0-certificate}
 (1+z)^2(3-2r)^2\frac{5\widetilde N_0}{24x^2}
   =a_0(r)+a_1(r)z+a_2(r)z^2,
\end{equation}
where
\begin{align*}
 a_0(r)&=(1-r)^3(3-2r)^2(10r^2+r+1),\\
 a_1(r)&=2(1-r)^3(3-2r)(9+5r+20r^2-4r^3),\\
 a_2(r)&=48(1-r)^3.
\end{align*}
All three coefficients are nonnegative for $-\frac12\leq r\leq1$.
For the first, the quadratic $10r^2+r+1$ has positive leading coefficient
and discriminant $1-40=-39<0$.  For the second,
\[
 9+5r+20r^2-4r^3=9+5r+r^2(20-4r)\geq\frac{13}{2}.
\]
The remaining factors in these expressions are visibly nonnegative.
Since the multiplier on the left of \eqref{rad:N0-certificate} is positive
on our domain, $N(0)\geq0$ follows.

The corresponding identity for $L$ is
\begin{equation}\label{rad:L-certificate}
 (1+z)^3(3-2r)^3\frac{25\widetilde L}{6x^2}
 =b_0(r)+b_1(r)z+b_2(r)z^2+b_3(r)z^3,
\end{equation}
where
\begin{align*}
 b_0(r)&=r^2(1-r)^3(3-2r)^3 f_0(r),\\
 b_1(r)&=(1-r)^3(3-2r)^2 f_1(r),\\
 b_2(r)&=2(1-r)^3(3-2r) f_2(r),\\
 b_3(r)&=24(1-r)^3(25-4r^2-10r),
\end{align*}
and every coefficient of the three remaining polynomials is as follows:
\begin{equation}\label{rad:f-polynomials}
\begin{aligned}
 f_0(r)&=480r^4-842r^3-117r^2+445r+160,\\
 f_1(r)&=336r^6-1008r^5+1096r^4-896r^3
                  +63r^2+685r+240,\\
 f_2(r)&=440r^4-1104r^3+368r^2+413r+210.
\end{aligned}
\end{equation}
Identities \eqref{rad:N0-certificate} and \eqref{rad:L-certificate}
follow by substituting \eqref{rad:z-substitution} into
\eqref{rad:N-expansion} and \eqref{rad:L-expansion}, respectively,
and clearing the displayed denominators.
The factor $25-4r^2-10r$ is at least $11$ on
$[-\frac12,1]$: its derivative is $-8r-10<0$ throughout that interval.
It remains to establish the signs of $f_0,f_1,f_2$.

Use the bijective change of variable
\[
 r=\frac{p-1/2}{1+p},\qquad
 p=\frac{r+1/2}{1-r}\geq0
 \quad\left(-\frac12\leq r<1\right).
\]
The resulting polynomial identities, with all coefficients displayed, are
\begin{equation}\label{rad:p-certificates}
\begin{aligned}
 4(1+p)^4 f_0\!\left(\frac{p-1/2}{1+p}\right)
    & =504p^4+4386p^3+12267p^2-1161p+174,\\
 4(1+p)^6 f_1\!\left(\frac{p-1/2}{1+p}\right)
    & =2064p^6+15486p^5+36669p^4+31368p^3\\
    &\qquad+20184p^2-2724p+522,\\
 4(1+p)^4 f_2\!\left(\frac{p-1/2}{1+p}\right)
    & =1308p^4+7650p^3+12366p^2-1842p+1044.
\end{aligned}
\end{equation}
Each right side is a polynomial with nonnegative coefficients in degrees
at least three, plus a strictly positive quadratic.  The exact
discriminants of those quadratic tails are
\begin{align*}
 (-1161)^2-4\cdot12267\cdot174&=-7\,189\,911<0,\\
 (-2724)^2-4\cdot20184\cdot522&=-34\,724\,016<0,\\
 (-1842)^2-4\cdot12366\cdot1044&=-48\,247\,452<0.
\end{align*}
Their leading coefficients are positive, so all three tails are strictly
positive on the real line.  Thus $f_i(r)>0$ for
$-\frac12\leq r<1$; their values at $r=1$ are, respectively,
$126,516,327$, also positive.  Every $b_j(r)$ is consequently nonnegative
on $[-\frac12,1]$.  Equation \eqref{rad:L-certificate} proves $L\geq0$,
finishing \eqref{rad:polynomial-goal} and hence
\eqref{rad:differential-bound}. The initial reduction now proves the lemma.
\end{proof}
